\chapter{From prototype to service}
\label{ch:operating-plan}

Picture the first ordinary Monday after the pilot. An author brings a private
technical document, an operator builds it, a reader tests it, and a sponsor
asks what the result is worth. Who owns each handoff, what does the customer
receive, and what happens when a parser fails? The twelve-week sequence gives us
a technical path. An operating service also needs clear boundaries, staffing,
procurement, privacy, maintenance, and an honest fallback if the full workflow
does not earn adoption.

The customer already pays experts to read technical documents. Humanvoice does
not need to invent a new reading market. It needs to reduce the repeated
diagnosis and clarification in research groups, investment teams, policy
institutions, and regulated technical operations. That keeps the commercial
question tied to the measured outcome.

\section{The work after the first build}
\label{sec:ops-work}

The feasibility build leaves four things that an operating service must make
routine:

\begin{enumerate}
\item a source adapter and build wrapper that can reproduce a document without
      a special developer's machine;
\item a finding and protected-object service that an author can use without
      exposing private source to an unapproved endpoint;
\item a pre-human release gate that catches ordinary integrity, argument,
      evidence, and pacing failures before a reader packet exists;
\item a reader study that produces interpretable observations rather than a
      style score; and
\item a decision packet that tells a sponsor whether to proceed, revise, or
      stop.
\end{enumerate}

The service should not promise continuous autonomous editing. A normal run is a
bounded authoring path for one document or one chapter: brief, blueprint, unit
draft, preflight, repair, release. The author chooses when to open a new
version, and the reader receives a frozen PDF only after the release gate
passes. This cadence makes the cost visible, keeps routine debugging out of the
reader's hands, and limits the blast radius of a parser or model failure.

\begin{table}[H]
\centering
\small
\begin{tabularx}{\textwidth}{@{}p{0.22\textwidth}p{0.23\textwidth}L@{}}
\toprule
Operating activity & Owner & Observable output \\
\midrule
Intake and authoring brief & Document owner & Decision, audience, named
exemplars, known vocabulary, exemption zones, disclosure, privacy boundary,
and source location \\
Blueprint and unit plan & Document owner with domain reviewer & Claim map,
dependencies, evidence anchors, section jobs, and visual plan \\
Snapshot and build & Technical operator & Reproducible PDF, build log, and
protected-object manifest \\
Pre-human preflight & Product operator & Gate results, mutation record,
abstentions, repair task, and release decision \\
Author review & Author with product support & Accepted, rejected, or deferred
findings, pacing and first-use dispositions, and a child version \\
Independent reading & Named domain reader & Timed answers, stopping point,
quoted span, missing relation, questions, and decision \\
Result interpretation & Sponsor and evaluator & Comparison, cost range, and
next action \\
Maintenance & Technical owner & Fixture replay, dependency update, and
incident record \\
\bottomrule
\end{tabularx}
\caption{A service run has a small number of human owners and visible outputs.}
\label{tab:ops-activities}
\end{table}

\section{Staffing and the twelve-week sequence}
\label{sec:ops-staffing}

The first team need not be large, but its responsibilities must be distinct.
The person implementing a parser should not be the only person deciding that a
reader understood the result. A practical starting team is a technical lead, a
document and data engineer, a domain editor, and a part-time evaluator or
research assistant. The sponsor supplies the decision and access to the named
readers; the team does not invent those roles.

The weeks are organized by the dependency of the work rather than by software
feature labels:

\begin{description}[style=nextline,leftmargin=31mm,labelindent=0mm]
\item[Weeks 1--2: establish the case] Select the live document, write the
reading brief, freeze the incumbent path, and extract protected objects by
hand for a small reference set. The output is a human-readable baseline and a
list of constructs the parser must handle.
\item[Weeks 3--4: make the source legible] Implement the snapshot and build
wrapper, source locations, bibliography identity checks, and a first detexed
view. Replay the positive and negative parser fixtures before adding prose
rules.
\item[Weeks 5--6: locate useful findings] Adapt a small set of Vale, textlint,
and project-specific rules, then run the local concept-density and first-use
prototype against the named exemplars. Add context windows and an author
decision form. Measure false positives, investigation time, and hit counts.
Keep a diagnostic only when the time it saves exceeds the time spent dismissing
it.
\item[Weeks 7--8: protect revisions] Implement object normalization, matching,
visual comparison, and unresolved correspondence handling. Run the deliberately
bad number, citation, equation, and macro fixtures.
\item[Weeks 9--10: run the reader path] Render incumbent and revised versions,
randomize order where possible, recruit the named readers, and record time,
answers, questions, and disclosure.
\item[Weeks 11--12: interpret and price] Rebuild from a clean environment,
replay the fixtures, calculate the predeclared outcomes and cost sensitivity,
and write the proceed, revise, or stop recommendation.
\end{description}

The sequence leaves room for a failed fixture. A parser that cannot preserve a
custom macro in week 4 should narrow the first release, not be hidden until
week 12. A reader who cannot answer the action question in week 10 should cause
an author revision or a stop, not a new dashboard.

\section{A small operating service}
\label{sec:ops-service}

An operator can run the first service as a local command-line application, with
an optional review page on the organization's network. The command line remains
the source of truth: it is scriptable and keeps private documents out of an
unexamined hosted system. The web page is only a convenience for viewing a
finding, accepting a revision, and presenting a clean reader packet.

The service-level promises should be concrete:

\begin{itemize}
\item A successful snapshot returns a PDF, source hash, build log, and
      protected-object manifest.
\item An unsupported source construct is reported with a location and does not
      receive a false "protected" result.
\item A reader packet contains no finding identifiers, private paths, or model
      confidence unless the document owner explicitly chooses to disclose them.
\item A revision with unresolved protected correspondence cannot be marked
      ready for a reader.
\item Every model suggestion records its endpoint, model version, prompt
      template, sampling parameters, and network policy.
\end{itemize}

These are operating behaviors, not claims that the product is safe in every
environment. The service owner must test them at installation and after each
dependency update. A failed promise creates an incident with a reproducible
source snapshot and a route back to the incumbent workflow.

\section{Procurement, licensing, and privacy}
\label{sec:ops-procurement}

The software inventory contains permissively licensed source projects, remote
services, and candidates whose license or maintenance status still needs
review. Procurement should treat each dependency as a decision with three
questions: may the organization use and modify it, can it be rebuilt from a
pinned source, and what data does it send outside the trust boundary?

A package can be useful without becoming a runtime dependency. For example, a
reference parser can generate expected spans for a fixture while the production
adapter remains a smaller implementation. A detector can supply an aggregate
research statistic while being forbidden from labeling an individual author.
The disposition should be written in terms of the job and the boundary, not in
terms of a package's popularity.

Privacy review follows the document path. The owner classifies source material
as public, internal, confidential, or restricted. The default local mode
allows all four classes subject to local access control. A remote model is
allowed only for public or explicitly redacted material, with a record of the
authorization and provider retention terms. Logs must not copy full source
passages by default; a finding stores the minimum context needed to reproduce
the decision, with access controlled separately from the reader PDF.

Figure~\ref{fig:trust-boundary} turns that policy into a data path. Any remote
call crosses an explicit approval point; the default manuscript path remains
inside the local zone.

\hvFigureTrustBoundary

The product should also support deletion and export. An organization must be
able to remove a source snapshot and its derived model prompts, or to export
the lineage needed for an authorized audit. A PDF hash is not a substitute for
a retention policy. These requirements are inexpensive to state at the start
and expensive to retrofit after a live document has crossed a boundary.

\section{Maintenance and quality of service}
\label{sec:ops-maintenance}

Writing tools change quickly. A package update can alter a parser tree, a model
update can alter suggestions, and a TeX package update can alter page breaks.
The service needs a small regression corpus with four kinds of expected result:

\begin{enumerate}
\item exact source and PDF invariants for protected objects;
\item expected findings and intentional abstentions for prose fixtures;
\item privacy and network checks for local and remote modes; and
\item reader-packet snapshots checked by a human after material changes.
\end{enumerate}

The first three can run in continuous integration. The fourth cannot be
reduced to an image hash because a page can remain pixel-identical while a
caption or cross-reference becomes confusing. A designated reader reviews the
changed pages and records whether the decision path remains clear.

Maintenance effort is part of the economic model. Let $u$ be annual document
volume, $m$ the minutes of operator maintenance per document, and $c_o$ the
loaded operator cost. The variable maintenance burden is
\begin{equation}
  C_{maint,var}=u\,m\,c_o/60.
  \label{eq:ops-maintenance}
\end{equation}
Fixed dependency review, security patches, and fixture refreshes add a fixed
term. A workflow that saves reader time but requires an operator to repair
every source manually may still be worthwhile for high-value documents, but
the buyer must see that trade.

\section{Three ways into the market}
\label{sec:ops-market}

Three likely customers have different operating needs.

\begin{description}[style=nextline,leftmargin=34mm,labelindent=0mm]
\item[Research groups] The buyer is a principal investigator or research
director; the author is a senior researcher; the reader is a collaborator,
reviewer, or funder. The service can begin as a guided review of a few high-
value proposals and papers.
\item[Investment and policy teams] The buyer pays for faster senior review of
model notes and memoranda. Confidentiality and citation provenance matter more
than broad style coverage, and the reader may have a short decision window.
\item[Regulated technical operations] The buyer needs a traceable revision
history and a controlled local deployment. The sales cycle is longer, but the
cost of an unexplained change is higher and the protected-object comparison may
be a direct procurement requirement.
\end{description}

The first customer should be chosen by document recurrence and reader cost, not
by the largest potential market. A single organization that can provide a
stable document class, a willing author, and independent readers is more useful
for learning than many one-off demos. The pilot should measure willingness to
repeat the workflow and willingness to pay for the reader outcome. A click on a
suggestion does not answer either question.

\section{A narrower product if the full service fails}
\label{sec:ops-alternatives}

The same components can support narrower offers. A "protected revision packet"
can combine source snapshots, object extraction, and visual comparison without
model suggestions. A "reader reconstruction study" can provide timed
evaluation for organizations that already have editing tools, using readers
unfamiliar with the project. A
"citation and equation safety service" can focus on high-risk technical
documents. These alternatives preserve the central value proposition and make
the result of a failed full-stack trial economically useful.

The reverse alternatives are less attractive. Selling a generic humanizer or
an authorship detector would put Humanvoice into crowded, weakly identified
markets and would conflict with the evidence on diversity, detector bias, and
citation reliability. A surface-only service may have a simpler sales story,
but it would not solve the reader problem documented in Part~I.

\section{Risk and contingency}
\label{sec:ops-contingency}

The principal operating risks are manageable when named early:

\begin{table}[H]
\centering
\small
\begin{tabularx}{\textwidth}{@{}p{0.23\textwidth}L p{0.29\textwidth}@{}}
\toprule
Risk & Early signal & Contingency \\
\midrule
Parser coverage remains low & More than a small fraction of protected objects
are unresolved & Narrow the supported LaTeX subset and provide a manual review
path \\
Readers see no improvement & Time and task completion do not move in the
feasibility case & Sell the measurement and protected-comparison service only,
or stop \\
Authors reject findings & Investigation burden exceeds the observed benefit &
Reduce rule families, change ordering, and test a diagnostic-only arm \\
Privacy review blocks inference & Source cannot be redacted without losing
meaning & Use deterministic local tools and postpone model suggestions \\
Maintenance dominates cost & Fixture failures follow dependency updates & Pin
versions, reduce the runtime stack, and price operator time explicitly \\
\bottomrule
\end{tabularx}
\caption{Contingencies preserve a useful service even when a larger product
claim fails.}
\label{tab:ops-risks}
\end{table}

\section{The operating decision}
\label{sec:ops-decision}

At scale, Humanvoice should be sold as a reduction in expensive reader work
with a visible meaning boundary. The buyer pays for a review path, not for a
promise that a model has made prose human. The author remains responsible for
the argument. The reader remains the authority for acceptance. The operator
keeps the source and software reproducible.

That arrangement is deliberately less glamorous than an autonomous editor,
but it has a measurable customer and a credible failure response. The
feasibility build should return three numbers that determine whether the
arrangement is worth extending: reader attention saved, protected-object error
rate, and total operating cost per document. Everything else is evidence for
understanding those numbers.
